\documentclass[conference]{IEEEtran}
\IEEEoverridecommandlockouts
\usepackage{ctex}
% 在无 Latin Modern 字体的环境中使用系统 CJK 字体覆盖拉丁字母回退字体。
\setmainfont{Noto Serif CJK SC}
\usepackage{cite}
\usepackage{amsmath,amssymb,amsfonts}
\usepackage{graphicx}
\usepackage{booktabs}
\usepackage{url}
\usepackage{xcolor}
\usepackage{hyperref}
\hypersetup{hidelinks}

\title{面向毫米级联合位姿到达的局部收敛增强软演员评论家方法}

\author{\IEEEauthorblockN{作者一\textsuperscript{*}}
\IEEEauthorblockA{\textit{所属院系，所属机构}\\
城市，国家\\
email@example.com}
\and
\IEEEauthorblockN{作者二}
\IEEEauthorblockA{\textit{所属院系，所属机构}\\
城市，国家\\
email@example.com}
}

\begin{document}
\maketitle

\begin{abstract}
深度强化学习已经能够解决随机末端位姿到达任务，但当成功容差从厘米级逐步收紧到毫米级时，策略常在目标邻域出现联合约束不对齐、修正方向不稳定以及确定性保持失败。本文研究在不缩小目标分布和不引入解析控制器的条件下，如何改善软演员评论家（Soft Actor-Critic，SAC）在高精度联合位姿控制中的局部收敛能力。我们首先构建 Hybrid KP-SAC-weak 表征基线，其观测同时包含显式六维位姿误差与已有三关键点联合几何误差，并保留稠密关键点 tracking/progress 奖励和精度课程；完整方法进一步从 Direction、Objective 和 Data 三个层面增强局部收敛：分别引入关键点局部 Jacobian，组合 precision quality、hold/leave 与终止门控瓶颈塑形，并组合成功轨迹回放与精度边界回放。训练难度按照位置与姿态容差从 L0 逐步收紧至 L14，使用固定的多种子确定性探针和确认评估控制升档。本文给出弱化表征基线、完整方法和三项关键子项消融的实验协议，在 L6、L9、L11、L12、L13 和 L14 分别检验从厘米级到严格姿态约束的联合位姿到达。主比较衡量三模块框架相对弱化基线的整体收益，关键子项的归因则来自 Full 与 M1--M3 的成对比较。
\end{abstract}

\begin{IEEEkeywords}
强化学习，软演员评论家，机械臂控制，联合位姿，课程学习，局部收敛
\end{IEEEkeywords}

\section{引言}
随机六维末端位姿到达要求策略同时处理三维位置和三维姿态。SAC 以最大熵目标提高了连续控制中的探索效率\cite{haarnoja2018sac}，课程学习则可以先训练较容易的任务，再逐步收紧目标容差\cite{bengio2009curriculum}。这些方法使机械臂能够在仿真中完成较宽松的随机位姿到达，但它们并不能保证策略在目标邻域继续产生方向正确、幅值合适且能够稳定保持的微小动作。

已有研究已经报告了随机位姿强化学习控制的可行性。例如，MAPPO-IK 在随机位置与姿态任务上获得了毫米级的位置误差和较小的姿态误差\cite{zhao2024mappoik}；连续位姿表示与课程学习也被用于缓解位姿表示不连续带来的训练问题\cite{hu2025rtpc}。然而，现有结果通常将“到达目标”和“在严格联合阈值内稳定保持”合并报告，对成功容差连续收紧后目标邻域的局部退化缺少系统分析。

本文关注以下问题：在目标分布、动作空间和训练预算保持不变时，如何降低高精度联合位姿控制的局部收敛退化？我们将退化归纳为三个相互作用的来源：局部状态不能充分区分可行修正方向，位置和姿态的优化信号发生竞争，以及普通经验回放中决定最终收敛的精度边界样本不足。围绕这三个来源，本文构建一个由局部方向信息、联合优化目标和精度边界数据组成的完整方法。

本文的主要贡献如下：
\begin{itemize}
  \item 提出一个从 L0 到 L14 的联合位姿精度课程，将研究重点从“是否能够到达”推进到“在毫米级容差下是否仍能局部收敛”；
  \item 在 Hybrid KP-SAC-weak 基线之上，从 Direction、Objective 和 Data 三个层面构建完整高精度方案，分别解决局部方向辨识、联合目标对齐与关键局部经验不足；
  \item 设计统一的三种子确定性探针、确认评估和冻结策略评估协议，使升档和论文结果不依赖训练日志中的在线随机成功率；
  \item 通过 Hybrid KP-SAC-weak 与 Full 检验三模块框架的整体收益，并通过 Full 与 M1--M3 成对比较 Jacobian、Bottleneck 和 Boundary Replay 三个关键子项的增量作用。
\end{itemize}

\section{相关工作}
\subsection{连续控制与课程学习}
SAC 通过最大熵强化学习目标在探索和回报之间取得平衡\cite{haarnoja2018sac}。课程学习从易到难组织训练任务，能够改善难任务的可学习性\cite{bengio2009curriculum}；面向精确运动规划的自动课程方法也表明，按任务复杂度组织训练分布有助于机械臂学习\cite{cho2025autocurriculum}。在机械臂场景中，课程通常通过逐步收紧位置和姿态容差来构造，但课程本身主要解决全局探索问题，不能直接解释目标邻域内的细粒度动作生成。

\subsection{位姿表示与经验回放}
六维位姿中的姿态表示会影响误差连续性和策略梯度。近期工作使用连续位姿表示、HER 和课程学习改善自由漂浮空间机械臂的轨迹规划\cite{hu2025rtpc}。HER 通过重新标记目标提高稀疏目标任务中的有效经验比例\cite{andrychowicz2017her}；优先经验回放则按样本重要性调整训练分布\cite{schaul2016sac}。Portela 等人提出的全身末端执行器位姿跟踪方法使用以末端为中心、随姿态旋转的立方体关键点，用三个具有直接对应关系的三维关键点同时表达位置和姿态，并报告了相较四元数、欧拉角和六维表示更易调节的联合跟踪效果\cite{portela2025wholebody}。近期的观测—动作空间对齐研究进一步表明，即使两种观测包含等量信息，和控制接口及奖励结构更匹配的表示也可能获得明显更高的到达成功率\cite{bhat2026alignment}。本文明确采用关键点位姿表达作为几何基础，但研究对象改为固定基座机械臂的随机位姿到达，并进一步加入 Jacobian 局部方向输入、联合瓶颈塑形和精度边界回放。因此，关键点表达本身属于已有方法基础，本文的增量在于面向毫米级强化学习到达任务的局部收敛增强与系统消融。

复杂奖励中的多个目标可能产生梯度竞争；近期工作因此研究了奖励课程以及位置、姿态等任务项的非线性融合\cite{freitag2025complexreward,jiang2025rfm}。这支持本文将位置—姿态竞争作为独立问题分析，而不是把所有奖励项简单相加。对于经验回放，聚类回放和分布感知采样分别从状态覆盖和数据分布保持角度说明了选择性采样的价值\cite{li2022clusteringreplay,garg2026distributionreplay}；本文将这一思想具体化为精度边界回放，并通过 M3 检验其作用。

\section{问题定义}
考虑一个六自由度机械臂。时刻 $t$ 的关节状态为 $q_t$，策略输出关节速度动作 $a_t=\dot q_t$，末端位姿为 $T(q_t)=[R(q_t),p(q_t)]$，目标位姿为 $T_g=[R_g,p_g]$。位置误差和姿态误差分别写为
\begin{equation}
e_p=\|p(q_t)-p_g\|_2,
\qquad
e_R=\|\log(R_g^{\mathsf T}R(q_t))\|_2 .
\label{eq:pose_error}
\end{equation}
给定位置阈值 $\epsilon_p^l$ 和姿态阈值 $\epsilon_R^l$，档位 $l$ 的联合成功事件为
\begin{equation}
\mathcal{S}_l=\mathbb{I}[e_p\leq\epsilon_p^l\;\wedge\;e_R\leq\epsilon_R^l].
\label{eq:success}
\end{equation}
训练目标是在完整随机目标分布上最大化 $\Pr(\mathcal{S}_l=1)$，而不是通过缩小目标距离或姿态差来获得表面上的高成功率。

\section{方法}
\subsection{整体框架与弱化表征基线}
图示中的方法增量统一归纳为三个复合模块：Direction 使用 Keypoint Jacobian；Objective 组合 precision quality、hold/leave 与 terminal-gated bottleneck shaping；Data 组合 Success Replay 与 Precision Boundary Replay。Hybrid KP-SAC-weak 同时保留显式位置—姿态误差与 Portela 式三关键点误差输入，并使用关键点距离/进度奖励、100\% ordinary current-level replay 和均衡目标采样；“Hybrid” 指两类观测并存，并不表示额外加入一套显式六维奖励。该基线不包含三个复合增强模块。Full 在同一环境、目标分布、网络规模、课程和预算下启用完整 Direction--Objective--Data 框架。因此二者的主比较衡量整体框架收益；M1、M2 分别删除 Objective/Data 中的 Bottleneck 和 Boundary Replay 子项，M3 删除完整 Direction 模块。

\subsection{关键点与局部运动学输入}
本文的关键点构造直接参考 Portela 等人的关键点位姿表示\cite{portela2025wholebody}：在末端坐标系中构造以末端位置为中心、方向随末端姿态旋转的立方体，并选取三个具有直接对应关系的顶点。三个非共线三维点足以同时确定平移和旋转；与直接拼接欧拉角或四元数相比，这种几何表示把联合位姿误差转换为同一坐标空间中的点误差。针对本文的固定基座 UR5，立方体边长和关键点位置按照工具坐标系定义，并随目标位姿同步生成。设末端及工具上选取的关键点坐标为 $x_i(q)$，目标关键点为 $x_i^g$。关键点距离为
\begin{equation}
d_k=\frac{1}{K}\sum_{i=1}^{K}\|x_i(q)-x_i^g\|_2 .
\label{eq:keypoint}
\end{equation}
对每个关键点，局部速度满足 $\dot x_i=J_i(q)\dot q$，其中 $J_i$ 为关键点 Jacobian。近期 Jacobian 探索式强化学习工作也采用局部 Jacobian 特征增强状态表示，以改善复杂机器人中的局部可观测性和探索效率\cite{tian2025jedprl}。Full 将归一化的关键点误差、关键点进度和局部 Jacobian 信息与标准位姿误差拼接为策略输入。该输入提供目标邻域的可辨识局部方向；M3 去除 Jacobian 输入，用于检验此信息是否带来额外收益。与该工作不同，本文不额外执行用于系统辨识的探索阶段，而是直接使用仿真中的解析关键点 Jacobian 作为策略输入，并把作用限定为高精度目标邻域的动作方向辨识。

\subsection{联合优化信号}
为避免位置或姿态单独变好而掩盖另一分量，Full 使用位置、姿态和关键点进度的连续误差下降信号。令 $\Delta e_p=e_{p,t-1}-e_{p,t}$、$\Delta e_R=e_{R,t-1}-e_{R,t}$，一种简化表达为
\begin{equation}
r_t=r_{\mathrm{task}}+\lambda_p\Delta e_p+\lambda_R\Delta e_R
 +\lambda_k\Delta d_k-r_{\mathrm{cost}} .
\label{eq:reward}
\end{equation}
其中 $r_{\mathrm{task}}$ 表示联合成功奖励，$r_{\mathrm{cost}}$ 包含动作和碰撞代价。Hybrid KP-SAC-weak 只保留关键点跟踪、关键点进度、连续成功奖励及必要代价；Full 与 M1--M3 额外共享按当前档位归一化的 precision quality 和 hold/leave。Full 的 joint bottleneck shaping 在这些阈值型信号之上，为最差归一化约束增加连续非终止势函数：
\begin{equation}
\Phi_b(s)=5\exp[-\rho_b(s)/2],\quad
r_{b,t}=(1-d_t)[\gamma\Phi_b(s_{t+1})-\Phi_b(s_t)].
\label{eq:bottleneck}
\end{equation}
其中 $d_t=1$ 表示 success、timeout、collision 或 hard failure terminal。所有 terminal transition 的 Bottleneck shaping 均为零，终止结果由专门的成功或失败项处理，因而不会产生与目标接近程度相关的隐藏终止奖惩。M1 去除该 bottleneck 项，但保留 Objective 的 precision quality 和 hold/leave。

\subsection{精度边界经验回放}
宽松课程阶段产生的经验不能充分覆盖严格阈值附近的状态。对 transition $(s_t,a_t,r_t,s_{t+1})$，Full 仅依据当前状态的 $0.8\leq\rho_{\max,t}\leq1.5$ 判定 Boundary subset；$s_{t+1}$ 只用于区分动作造成的保持、退出或单项失败。该定义直接表达“当前处于精度边界时应采取什么动作”，并与 Jacobian 方向分析和局部反向修正率的当前状态门控一致。该机制的目的不是简单重复失败样本，而是增加对“已经接近目标但还需联合修正”的样本访问。M2 去除 precision boundary replay，用于验证目标边界数据分布的作用。

\subsection{自动精度课程}
课程档位及阈值见表~\ref{tab:levels}。所有方法共享同一 precision action scaling：令 $z_t=\max(e_{p,t}/\epsilon_p,e_{R,t}/\epsilon_R)$，当 $z_t\leq1$ 时 $g(z_t)=0.05+0.05z_t$，否则 $g(z_t)=0.10+0.90\min(z_t-1,1)^2$，实际关节速度为 $\dot q_{\rm cmd}=g(z_t)(0.7a_t)$。因此跨档结果应解释为统一精度控制合同下的 strictness--performance curve。每档至少累计 $50\,000$ 个 transition；每累计 $300\,000$ 个 transition 触发一次当前档位 deterministic probe。Probe 使用三个评估种子、每种子 300 个 episode；三种子平均成功率至少为 85\%，最低单种子成功率至少为 80\%；碰撞率、关节越界率和 timeout 按协议记录后，再使用每种子 1000 个 episode 的确认评估。只有确认通过才升档。训练日志中的在线 stochastic success 仅作为诊断，不改变档位。

\begin{table}[t]
\caption{精度课程档位}
\label{tab:levels}
\centering
\small
\begin{tabular}{c|rrrrrrrrrrrr}
\toprule
档位 $l$ &0&1&2&3&4&5&6&7&8&9&10&11\\
\midrule
$\epsilon_p$ / mm&100&80&64&51.2&50&25&10&5&2.5&2&1.5&1\\
$\epsilon_R$ / rad&.30&.24&.192&.1536&.10&.10&.10&.10&.10&.10&.10&.10\\
\bottomrule
\end{tabular}
\end{table}

\section{实验设置}
\subsection{环境与训练}
实验在 PyBullet UR5 无障碍环境中进行。目标为随机六维末端位姿，动作是关节速度，所有方法使用相同的 episode 长度、SAC 网络、目标分布和 3,000,000 transition 预算。训练种子为 11001 和 11002；评估种子为 31001、31002 和 31003。每个 checkpoint 在冻结 Actor 后进行 deterministic 评估，主表报告 L6、L9、L11、L12、L13 和 L14。

\subsection{方法矩阵}
\begin{table}[t]
\caption{正式方法与验证问题}
\label{tab:methods}
\centering
\small
\begin{tabular}{p{0.22\columnwidth}|p{0.69\columnwidth}}
\toprule
方法 & 验证内容\\
\midrule
Hybrid KP-SAC-weak & 显式 6D 位姿误差、Portela 式关键点误差/奖励、100\% ordinary replay 和课程；测量弱化混合表征基线的退化边界。\\
Full & 加入 Jacobian、bottleneck 和 precision boundary replay；检验完整设计。\\
M1 & 去除 joint bottleneck shaping；检验联合优化信号。\\
M2 & 去除 precision boundary replay；检验目标边界样本。\\
M3 & 去除 Jacobian 输入；检验局部运动学方向信息。\\
\bottomrule
\end{tabular}
\end{table}

\subsection{评价指标}
除联合成功率外，报告碰撞率、timeout 率、关节越界率、终端位置和姿态误差的均值与 95 分位数，并按距离桶、姿态差桶和任务 cell 统计覆盖情况，主报告 cell P10 与最差 10 个 cell 的平均成功率，absolute worst cell 放入附录。进一步区分 position-only、orientation-only 和 joint-pose 失败，以判断性能差异来自位置、姿态还是联合瓶颈。

\section{结果与分析}
\subsection{Matched-level 辅助分析}
固定预算自动课程是主结果，因为更早通过课程并进入更高精度档位本身就是方法能力的一部分。为区分整体方案收益与高精度训练暴露量，在 Full 和 Hybrid KP-SAC-weak 都进入过的共同档位（优先 L6、L9）选择同档位、相近局部训练暴露量的 checkpoint，使用相同三个评估种子和每种子 1000 个 episode 进行冻结评估。主表报告 fixed-budget 结果，辅助表报告 matched-level 结果。前者衡量完整方法最终能推进到何种精度，后者检验在相同精度档位时 Full 是否已经具有局部收敛优势。

\subsection{主结果}
表~\ref{tab:mainresult} 是论文主结果的预留表。所有单元格在新实验完成后填入三个评估种子的均值和标准差；在此之前不使用旧 checkpoint 或旧评估 JSON。

\begin{table}[t]
\caption{不同精度档位的冻结确定性成功率（待填）}
\label{tab:mainresult}
\centering
\small
\begin{tabular}{c|cccccc|cccccc}
\toprule
&\multicolumn{6}{c|}{成功率均值 / \%}&\multicolumn{6}{c}{位置误差 P95 / mm}\\
方法 & L6 & L9 & L11 & L12 & L13 & L14 & L6 & L9 & L11 & L12 & L13 & L14\\
\midrule
Hybrid KP-SAC-weak & -- & -- & -- & -- & -- & -- & -- & -- & -- & -- & -- & --\\
Full & -- & -- & -- & -- & -- & -- & -- & -- & -- & -- & -- & --\\
M1 & -- & -- & -- & -- & -- & -- & -- & -- & -- & -- & -- & --\\
M2 & -- & -- & -- & -- & -- & -- & -- & -- & -- & -- & -- & --\\
M3 & -- & -- & -- & -- & -- & -- & -- & -- & -- & -- & -- & --\\
\bottomrule
\end{tabular}
\end{table}

预期分析顺序是先比较 Full 与 Hybrid KP-SAC-weak 在 L6、L9、L11、L12、L13 和 L14 的整体差距，再检查差距是否随容差收紧而扩大。如果 Full 只在宽松档位提高而在 L9--L14 没有改善，则不能声称改善了高精度局部收敛。三个复合模块的整体收益来自 Full 与弱基线；关键子项的作用依据 Full 与 M1--M3 的成对消融及对应退化模式，其中 M1/M2 分别删除 Objective/Data 中的一个子项，M3 删除完整 Direction 模块。

\subsection{消融结果}
\begin{table}[t]
\caption{机制消融结果（待填）}
\label{tab:ablation}
\centering
\small
\begin{tabular}{c|ccc|p{0.34\columnwidth}}
\toprule
方法 & L6 & L9 & L11 & 主要观察\\
\midrule
Full & -- & -- & -- & --\\
M1 & -- & -- & -- & 位置—姿态竞争与失败类型待分析\\
M2 & -- & -- & -- & 边界样本覆盖待分析\\
M3 & -- & -- & -- & 局部方向信息待分析\\
\bottomrule
\end{tabular}
\end{table}

\section{讨论与局限}
本文结论限定在无障碍 PyBullet 仿真和给定随机目标分布内。仿真没有真实编码器噪声、减速器间隙、标定误差、视觉误差和通信延迟，因此一毫米阈值应表述为仿真中的联合成功判据，而不是未经验证的真实硬件精度。后续工作需要加入动力学扰动、观测噪声和真实机器人验证。此外，M1/M2 只分别移除 Objective/Data 中的 Bottleneck 与 Boundary Replay 子项，因此它们支持这两个子项的增量归因，不等价于整模块删除；M3 才是完整 Direction 模块消融。

\section{结论}
本文提出了一个面向毫米级联合位姿到达的局部收敛增强 SAC 研究框架。Hybrid KP-SAC-weak 提供显式六维误差、三关键点联合误差、稠密关键点奖励和精度课程，用于刻画不含额外局部高精度辅助机制时的能力边界。通过弱化基线与完整方法的整体比较、Full 家族内部单模块消融和固定多种子评估，本文将检验完整方案能否把高精度退化边界推进到更严格的联合位姿约束，并分别分析局部方向、联合优化目标和精度边界数据的作用。实验结果完成后，将根据表~\ref{tab:mainresult} 和表~\ref{tab:ablation} 填入定量结论。

\section*{致谢}
本节在投稿前填写基金、平台和计算资源信息。

\begin{thebibliography}{00}
\bibitem{haarnoja2018sac} T. Haarnoja, A. Zhou, P. Abbeel, and S. Levine, ``Soft Actor-Critic: Off-Policy Maximum Entropy Deep Reinforcement Learning with a Stochastic Actor,'' in \emph{Proc. ICML}, vol. 80, pp. 1861--1870, 2018. [Online]. Available: \url{https://proceedings.mlr.press/v80/haarnoja18b.html}
\bibitem{bengio2009curriculum} Y. Bengio, J. Louradour, R. Collobert, and J. Weston, ``Curriculum learning,'' in \emph{Proc. ICML}, pp. 41--48, 2009, doi: 10.1145/1553374.1553380.
\bibitem{andrychowicz2017her} M. Andrychowicz \emph{et al.}, ``Hindsight experience replay,'' in \emph{Advances in Neural Information Processing Systems}, vol. 30, 2017. [Online]. Available: \url{https://papers.nips.cc/paper/7090-hindsight-experience-replay}
\bibitem{schaul2016sac} T. Schaul, J. Quan, I. Antonoglou, and D. Silver, ``Prioritized experience replay,'' in \emph{Proc. ICLR}, 2016. [Online]. Available: \url{https://arxiv.org/abs/1511.05952}
\bibitem{zhao2024mappoik} Y. Zhao \emph{et al.}, ``Inverse kinematics solution and control method of 6-degree-of-freedom manipulator based on deep reinforcement learning,'' \emph{Scientific Reports}, vol. 14, Art. no. 12467, 2024, doi: 10.1038/s41598-024-62948-6.
\bibitem{hu2025rtpc} X. Hu \emph{et al.}, ``Deep reinforcement learning-based trajectory planning with continuous pose representation for 6-DoF free-floating space robot,'' \emph{Aerospace Science and Technology}, Art. no. 110540, 2025, doi: 10.1016/j.ast.2025.110540.
\bibitem{portela2025wholebody} T. Portela, A. Cramariuc, M. Mittal, and M. Hutter, ``Whole-body end-effector pose tracking,'' arXiv:2409.16048v2 [cs.RO], 2025. [Online]. Available: \url{https://arxiv.org/abs/2409.16048}
\bibitem{tian2025jedprl} Y. Tian, C. K. Ng, and H. Ren, ``Jacobian Exploratory Dual-Phase Reinforcement Learning for Dynamic Endoluminal Navigation of Deformable Continuum Robots,'' arXiv:2509.00329v1 [cs.RO], 2025. [Online]. Available: \url{https://arxiv.org/abs/2509.00329}
\bibitem{cho2025autocurriculum} D.-S. Cho, J.-M. Cho, and W.-T. Kim, ``Guided deep reinforcement learning framework using automated curriculum scheme for accurate motion planning,'' \emph{Engineering Applications of Artificial Intelligence}, vol. 139, Art. no. 109541, 2025, doi: 10.1016/j.engappai.2024.109541.
\bibitem{bhat2026alignment} V. Bhat, Z. Suleymanova, and C. Bellinger, ``Beyond Information Sufficiency: Observation-Action Space Alignment in Robotic Reinforcement Learning,'' in \emph{Proc. Canadian Conf. Artificial Intelligence}, PMLR, vol. 318, pp. 1052--1059, 2026. [Online]. Available: \url{https://proceedings.mlr.press/v318/bhat26a.html}
\bibitem{freitag2025complexreward} K. Freitag, K. Ceder, R. Laezza, K. Åkesson, and M. H. Chehreghani, ``Curriculum Reinforcement Learning for Complex Reward Functions,'' arXiv:2410.16790v2 [cs.LG], 2025. [Online]. Available: \url{https://arxiv.org/abs/2410.16790}
\bibitem{jiang2025rfm} K. Jiang, Z. Fu, J. Guo, W. Zhang, and H. Chen, ``Learning Whole-Body Loco-Manipulation for Omni-Directional Task Space Pose Tracking With a Wheeled-Quadrupedal-Manipulator,'' \emph{IEEE Robotics and Automation Letters}, vol. 10, no. 2, 2025, doi: 10.1109/LRA.2024.3519856.
\bibitem{li2022clusteringreplay} M. Li, T. Huang, and W. Zhu, ``Clustering experience replay for the effective exploitation in reinforcement learning,'' \emph{Pattern Recognition}, vol. 131, Art. no. 108875, 2022, doi: 10.1016/j.patcog.2022.108875.
\bibitem{garg2026distributionreplay} A. Garg, A. R. Sharma, A. Arora, and S. Verma, ``Distribution-aware sampling of replay buffer for mitigating catastrophic forgetting,'' \emph{Information Sciences}, vol. 753, Art. no. 123640, 2026, doi: 10.1016/j.ins.2026.123640.
\bibitem{astorquia2025benchmark} M. Astorquia \emph{et al.}, ``Comparative benchmark of sampling-based and DRL motion planning methods for industrial robotic arms,'' \emph{Sensors}, vol. 25, no. 17, Art. no. 5282, 2025, doi: 10.3390/s25175282.
\end{thebibliography}

\end{document}
